\documentclass[11pt, a4paper]{article}

% ── Codificación y Lenguaje ──────────────────────────────────────────────────
\usepackage[utf8]{inputenc}
\usepackage[T1]{fontenc}
\usepackage[spanish, es-tabla]{babel}

% ── Geometría ───────────────────────────────────────────────────────────────
\usepackage[top=2.5cm, bottom=2.5cm, left=2.5cm, right=2.5cm, headheight=24pt]{geometry}

% ── Matemáticas ─────────────────────────────────────────────────────────────
\usepackage{amsmath, amssymb, amsthm}
\usepackage{mathtools}

% ── Colores y Cajas ─────────────────────────────────────────────────────────
\usepackage{xcolor}
\usepackage[most]{tcolorbox}
\tcbuselibrary{skins, breakable, theorems}

% ── Tablas ──────────────────────────────────────────────────────────────────
\usepackage{booktabs}
\usepackage{array}
\usepackage{multirow}
\usepackage{longtable}

% ── Listas ──────────────────────────────────────────────────────────────────
\usepackage{enumitem}

% ── Encabezado y Pie de Página ──────────────────────────────────────────────
\usepackage{fancyhdr}

% ── Secciones ───────────────────────────────────────────────────────────────
\usepackage{titlesec}

% ── Gráficos y Diagramas ────────────────────────────────────────────────────
\usepackage{tikz}
\usetikzlibrary{shapes.geometric, arrows.meta, positioning, calc}
\usepackage{graphicx}

% ── Columnas múltiples ──────────────────────────────────────────────────────
\usepackage{multicol}

% ── Espaciado ───────────────────────────────────────────────────────────────
\usepackage{setspace}
\setlength{\parskip}{4pt}
\setlength{\parindent}{0pt}

% ── Paleta GOP ──────────────────────────────────────────────────────────────
\definecolor{gopblue}{RGB}{31, 78, 121}
\definecolor{gopcyan}{RGB}{70, 130, 180}
\definecolor{gopgreen}{RGB}{0, 100, 0}
\definecolor{gopgreenlight}{RGB}{230, 255, 230}
\definecolor{goporange}{RGB}{200, 100, 0}
\definecolor{goporangelight}{RGB}{255, 245, 210}
\definecolor{gopred}{RGB}{160, 0, 0}
\definecolor{gopredlight}{RGB}{255, 230, 230}
\definecolor{gopgray}{RGB}{90, 90, 90}
\definecolor{gopgraylight}{RGB}{245, 245, 240}
\definecolor{gopbluedef}{RGB}{0, 70, 140}
\definecolor{gopbluedeflight}{RGB}{220, 235, 255}

% ── Hipervínculos y TOC ─────────────────────────────────────────────────────
\usepackage[colorlinks=true, linkcolor=gopblue, urlcolor=gopblue]{hyperref}

% ── Entornos tcolorbox ───────────────────────────────────────────────────────
\newtcolorbox{definicion}[1]{
  enhanced, breakable,
  colback=gopbluedeflight, colframe=gopbluedef,
  fonttitle=\bfseries\small, title={Definición: #1}, coltitle=white,
  attach boxed title to top left={yshift=-2mm, xshift=6pt},
  boxed title style={colback=gopbluedef, colframe=gopbluedef, sharp corners},
  sharp corners=south, top=4mm, before skip=6pt, after skip=6pt,
}
\newtcolorbox{formulas}[1]{
  enhanced, breakable,
  colback=gopgreenlight, colframe=gopgreen,
  fonttitle=\bfseries\small, title={Fórmulas: #1}, coltitle=white,
  attach boxed title to top left={yshift=-2mm, xshift=6pt},
  boxed title style={colback=gopgreen, colframe=gopgreen, sharp corners},
  sharp corners=south, top=4mm, before skip=6pt, after skip=6pt,
}
\newtcolorbox{tipprueba}[1]{
  enhanced, breakable,
  colback=goporangelight, colframe=goporange,
  fonttitle=\bfseries\small, title={TIP PRUEBA: #1}, coltitle=white,
  attach boxed title to top left={yshift=-2mm, xshift=6pt},
  boxed title style={colback=goporange, colframe=goporange, sharp corners},
  sharp corners=south, top=4mm, before skip=6pt, after skip=6pt,
}
\newtcolorbox{trampa}[1]{
  enhanced, breakable,
  colback=gopredlight, colframe=gopred,
  fonttitle=\bfseries\small, title={TRAMPA/ADVERTENCIA: #1}, coltitle=white,
  attach boxed title to top left={yshift=-2mm, xshift=6pt},
  boxed title style={colback=gopred, colframe=gopred, sharp corners},
  sharp corners=south, top=4mm, before skip=6pt, after skip=6pt,
}
\newtcolorbox{ejercicio}[1]{
  enhanced, breakable,
  colback=gopgraylight, colframe=gopgray,
  fonttitle=\bfseries\small\color{black}, title={Ejercicio: #1},
  attach boxed title to top left={yshift=-2mm, xshift=6pt},
  boxed title style={colback=gopgray!40, colframe=gopgray, sharp corners},
  sharp corners=south, top=4mm, before skip=6pt, after skip=6pt,
}

% ── Encabezado y pie ─────────────────────────────────────────────────────────
\pagestyle{fancy}
\fancyhf{}
\fancyhead[L]{\small\textbf{GOP ICS3213} --- Compendio de Ejercicios: Variabilidad y Teoría de Colas}
\fancyhead[C]{}
\fancyhead[R]{\small\thepage}
\fancyfoot[L]{\footnotesize Solución Oficial --- Transcripción en LaTeX}
\fancyfoot[R]{\small\thepage}
\renewcommand{\headrulewidth}{0.4pt}
\renewcommand{\footrulewidth}{0.4pt}

% ── Estilo de secciones ──────────────────────────────────────────────────────
\titleformat{\section}
  {\color{gopblue}\Large\bfseries}{\color{gopblue}\thesection.}{0.5em}{}
  [\color{gopblue}\titlerule]
\titleformat{\subsection}
  {\color{gopblue}\normalsize\bfseries}{\color{gopblue}\thesubsection.}{0.5em}{}
\titleformat{\subsubsection}
  {\color{gopblue}\small\bfseries}{\color{gopblue}\thesubsubsection.}{0.5em}{}

\begin{document}

% ════════════════════════════════════════════════════════════════════════════
% PORTADA
% ════════════════════════════════════════════════════════════════════════════
\begin{titlepage}
  \centering
  \vspace*{2cm}
  {\Huge\bfseries\color{gopcyan} Compendio de Ejercicios\par}
  \vspace{0.4cm}
  {\LARGE\bfseries Variabilidad y Teoría de Colas\par}
  \vspace{0.3cm}
  {\large Gestión de Operaciones (ICS3213)\par}
  \vspace{1.2cm}
  \rule{\textwidth}{0.4pt}
  \vspace{0.4cm}
  {\small Material extraído de pautas oficiales de Interrogaciones I2.\par}
  \vspace{0.4cm}
  \rule{\textwidth}{0.4pt}
  \vfill
\end{titlepage}

\newpage
\section*{Variabilidad y Teoría de Colas}

\subsection*{[I2_2022] Pregunta 2: Gestión de Capacidad y Variabilidad (Teoría de Colas) (20 Puntos)}
\begin{tipprueba}{¿Cuándo usar este enfoque?}
    \textbf{Identificación:} Se analiza el diseño de capacidad y la variabilidad de una cola usando el balance de costos operativos y de espera de los clientes, seguido de un análisis en cadena (tándem) donde se propaga la variabilidad.\\
    \textbf{Qué aplicar:}
    \begin{itemize}
      \item En M/M/1, el tiempo de permanencia es $W_s = \frac{1}{\mu - \lambda}$.
      \item El costo total es $CT(\mu) = CE(W_s) \cdot \text{clientes} + 10\mu$. La pauta utiliza una minimización simplificada por cliente individual: $CT(\mu) = 10 + \frac{1000}{\mu - 25} + 10\mu$.
      \item En sistemas en tándem, el coeficiente de variación de salida de la estación 1 es $C_{d,1}^2 \approx \rho_1^2 \cdot C_{s,1}^2 + (1-\rho_1)^2 \cdot C_{a,1}^2$, que alimenta como $C_{a,2}^2$ a la estación 2 (G/G/1).
    \end{itemize}
  \end{tipprueba}

  \textbf{Enunciado:}
  Usted está pensando en colocar un emprendimiento de comida y debe determinar la capacidad de atención. Para llevar esto a cabo, determina el costo total de espera $CE(W_s)$ de los clientes en el sistema, el cual depende del tiempo de permanencia ($W_s$) según la ecuación:
  \begin{equation*}
    CE(W_s) = 10 + 1000 W_s
  \end{equation*}

  La tasa de llegada de los clientes al puesto es de $\lambda = 25$ clientes por hora, con un coeficiente de variación de llegada $C_a = 1$. 

  La capacidad de atender a los clientes ($\mu$) se puede ajustar linealmente y tiene un costo operativo de $\$10/\text{unidad de capacidad}$. Es decir, si se contrata una tasa de servicio de $\mu$ clientes por hora, el costo del servicio es de $10\mu$. El coeficiente de variación de la atención es $C_s = 1$.

  \begin{itemize}
    \item[\textbf{a)}] (12 Puntos) Con esta información determine la capacidad óptima de atención ($\mu$) que debe implementar en su emprendimiento. (*HINT: Un sistema con $C_a=1$ y $C_s=1$ se comporta como un modelo M/M/1*).
    \item[\textbf{b)}] (8 Puntos) Usted ha comprado la capacidad productiva determinada en (a) pero olvidó que debía instalar una caja registradora antes del mesón de atención. Si la caja tiene una capacidad de atender $50$ clientes por hora y un coeficiente de variación de servicio $C_{s,\text{caja}} = 2$. ¿Cómo cambia el tiempo de espera desde que el cliente sale de la caja y es atendido en el mesón? ¿Aumenta o disminuye? Justifique mediante el análisis de propagación de variabilidad.
  \end{itemize}

  \medskip
  \begin{ejercicio}{Solución}
  \textbf{Desarrollo de la Solución:}

  \subsubsection*{a) Determinación de la Capacidad Óptima $\mu$}
  El sistema sin caja se comporta como una cola \textbf{M/M/1}.
  El tiempo en el sistema es $W_s = \frac{1}{\mu - 25}$.
  Minimizamos la función de costo simplificada por cliente (según pauta oficial):
  \begin{equation*}
    \min_{\mu} \quad CT = 10 + \frac{1000}{\mu - 25} + 10 \mu
  \end{equation*}

  Derivamos con respecto a $\mu$ e igualamos a cero:
  \begin{equation*}
    \frac{dCT}{d\mu} = -\frac{1000}{(\mu - 25)^2} + 10 = 0 \implies (\mu - 25)^2 = 100
  \end{equation*}
  \begin{equation*}
    \mu - 25 = 10 \implies \mu = \mathbf{35 \text{ clientes/hora}}
  \end{equation*}
  La capacidad óptima de atención recomendada es de \textbf{35 clientes/hora}.

  \subsubsection*{b) Impacto de la Caja Registradora (Propagación de Variabilidad)}
  Con la caja antes del mesón, se modela un sistema de colas en tándem:
  \begin{itemize}
    \item \textbf{Caja (Estación 1):} $\mu_1 = 50$, $C_{s,1} = 2$, $\lambda = 25 \implies \rho_1 = \frac{25}{50} = 0.5$.
    \item \textbf{Mesón (Estación 2):} $\mu_2 = 35$, $C_{s,2} = 1$, $\lambda = 25 \implies \rho_2 = \frac{25}{35} = 0.7143$.
  \end{itemize}

  El coeficiente de variación de las salidas de la caja ($C_{d,1}^2$) determina el coeficiente de variación de llegadas en el mesón ($C_{a,2}^2$):
  \begin{equation*}
    C_{a,2}^2 \approx C_{d,1}^2 \approx \rho_1^2 \cdot C_{s,1}^2 + (1 - \rho_1)^2 \cdot C_{a,1}^2
  \end{equation*}
  \begin{equation*}
    C_{a,2}^2 \approx (0.5)^2 \cdot (2)^2 + (1 - 0.5)^2 \cdot (1)^2 = 0.25 \cdot 4 + 0.25 \cdot 1 = 1.25
  \end{equation*}

  El mesón pasa a modelarse como una cola \textbf{G/G/1}. El tiempo de espera en la cola ($W_{q,2}$) se estima con la ecuación de Kingman (VUT):
  \begin{equation*}
    W_{q,2} = \left( \frac{C_{a,2}^2 + C_{s,2}^2}{2} \right) \left( \frac{\rho_2}{1 - \rho_2} \right) \left( \frac{1}{\mu_2} \right)
  \end{equation*}
  \begin{equation*}
    W_{q,2} = \left( \frac{1.25 + 1}{2} \right) \left( \frac{0.7143}{1 - 0.7143} \right) \left( \frac{1}{35} \right) \approx \mathbf{0.0804 \text{ horas}} \quad (4.82 \text{ min})
  \end{equation*}

  En la situación original sin caja (llegada directa con $C_{a}^2 = 1$):
  \begin{equation*}
    W_{q,\text{original}} = \left( \frac{1 + 1}{2} \right) \left( \frac{0.7143}{1 - 0.7143} \right) \left( \frac{1}{35} \right) \approx \mathbf{0.0714 \text{ horas}} \quad (4.29 \text{ min})
  \end{equation*}

  \textbf{Análisis:} El tiempo de espera en el mesón \textbf{aumenta} de 0.0714 horas a 0.0804 horas (un incremento del $12.6\%$). Esto se debe a que la caja introduce variabilidad en el flujo de clientes ($C_{s,\text{caja}} = 2$), la cual se propaga aguas abajo, congestionando el mesón y aumentando la demora promedio en la fila.

\end{ejercicio}

\newpage

\end{document}
